% Capitulo: Master data: de projeto eterno a oportunidade
\chapter{Master data: de projeto eterno a oportunidade}
\label{cap:master-data}

Toda empresa tem algumas entidades que aparecem em quase todos os processos: clientes,
produtos, fornecedores, locais, funcionários. Os dados que descrevem essas entidades são
chamados de dados mestres, ou master data. Quando eles estão bem cuidados, quase ninguém
repara. Quando não estão, tudo dá errado: a receita por cliente não bate, o mesmo produto
aparece com três códigos, o relatório de compras não conversa com o de estoque.

Por muito tempo, cuidar disso foi visto como um projeto caro, lento e difícil de
justificar, e muitas empresas preferiram conviver com as planilhas. Este capítulo
argumenta que essa conta mudou. Com a abordagem certa, e com a ajuda da IA, master data
pode deixar de ser um projeto eterno e virar uma oportunidade.

\section{O que é master data}
\label{sec:o-que-e-master-data}

Dados mestres são as entidades centrais do negócio, compartilhadas por vários processos e
sistemas \parencite{dama2017,loshin2009}. Eles se distinguem de dois outros tipos de dado.
O dado transacional registra eventos: a venda, a viagem, o lançamento contábil. Ele é
volumoso, cresce todo dia e quase nunca muda depois de registrado. O dado de referência
são listas de valores usadas para classificar outros dados: estados, moedas, códigos de
unidade de medida. Ele é pequeno, muda pouco e muitas vezes vem de fora da empresa.

O dado mestre fica no meio. É menos volumoso que o transacional, mas muda com frequência
(o cliente muda de endereço, o produto ganha uma nova embalagem) e é usado por quase todas
as transações. E, principalmente, costuma existir em vários sistemas ao mesmo tempo. O
ERP tem o seu cadastro de clientes, o CRM tem outro, a área comercial mantém uma planilha
com um terceiro. Cada um com os seus códigos e a sua forma de escrever os nomes
(Tabela~\ref{tab:registro-mestre}).

\begin{table}[htbp]
  \centering
  \caption{O mesmo cliente em três fontes e o registro mestre que as concilia.}
  \label{tab:registro-mestre}
  \small
  \begin{tabular}{@{}lll@{}}
    \toprule
    Fonte & ID na fonte & Nome como está na fonte \\
    \midrule
    ERP                & 4471          & AGRO SILVA LTDA \\
    CRM                & C-0932        & Agro Silva \\
    Planilha comercial & (sem ID)      & Agrosilva Faz. Boa Vista \\
    \midrule
    \textbf{Registro mestre} & \textbf{1001} & \textbf{Agro Silva Ltda} \\
    \bottomrule
  \end{tabular}
  \fonte{exemplo fictício elaborado pelo autor.}
\end{table}

Gerir master data é, no essencial, manter duas coisas. A primeira é o registro mestre
(\emph{golden record}): a versão única e confiável de cada entidade, com um identificador
próprio. A segunda é a tabela de correspondência, que liga o identificador de cada fonte ao
identificador mestre. Com as duas, qualquer análise pode juntar dados de sistemas
diferentes sabendo que está falando do mesmo cliente.

Quem leu o capítulo~\ref{cap:dw} já viu esse problema com outro nome. As dimensões
conformadas de Kimball, compartilhadas por todos os processos de negócio, são master data
aplicado ao data warehouse (\veja{sec:kimball}). Sem um cadastro mestre por trás, não há
dimensão conformada que se sustente.

\section{Como master data era visto}
\label{sec:master-data-antes}

O jeito tradicional de resolver master data era um programa grande, separado de qualquer
entrega de negócio. Primeiro se mapeavam todas as fontes de todos os domínios. Depois se
conciliavam os nomes e os códigos, registro por registro. Depois se definia a governança:
quem cria, quem altera, quem aprova. Só então havia alguma coisa para entregar.

Esse caminho exigia gente e tempo. Era preciso contratar analistas de dados para
conversar com os analistas de negócio de cada área, entender como cada sistema registrava
os dados, negociar regras que nenhuma área queria ceder. Até que o resultado aparecesse,
passavam meses, muitas vezes anos (parte de cima da Figura~\ref{fig:master-data-abordagens}).
E o resultado, quando aparecia, era difícil de mostrar: \enquote{dados mais consistentes}
não é algo que a diretoria vê numa reunião.

Diante disso, muitas empresas fizeram uma escolha racional: continuar com as planilhas. A
planilha com o \enquote{de-para} mantido à mão já existia, funcionava mais ou menos e não
pedia orçamento. O custo dela era real, mas espalhado: retrabalho a cada fechamento, cada
área com a sua versão, números que não batiam e consumiam a primeira meia hora de cada
reunião. Ninguém via esse custo numa linha do orçamento, e por isso ele nunca competia com
o custo visível de um programa de master data.

\rascunho{
  Relato: uma situação em que a falta de master data causou um problema concreto
  (números divergentes, retrabalho no fechamento, decisão tomada com base errada). Sem
  nomes e com valores só em ordem de grandeza.
}

\section{Master data puxado por um projeto de negócio}
\label{sec:master-data-puxado}

A abordagem que tem funcionado melhor na minha experiência inverte a lógica. Em vez de
começar um programa de master data e procurar quem vai usá-lo, começa-se por um projeto de
negócio que já foi aprovado, já tem patrocinador e orçamento, e identificam-se os braços
desse projeto que dependem de master data.

Um exemplo típico: um projeto comercial precisa juntar a carteira de clientes de três
sistemas para acompanhar o desempenho de cada vendedor. Sem um cadastro mestre de clientes,
o projeto não fica de pé, porque o mesmo cliente aparece três vezes com nomes diferentes. É
esse o momento de abrir, em paralelo, um projeto de master data só para o domínio de
clientes. A necessidade já está justificada, porque sem ele o projeto de negócio não
entrega. O recurso já está garantido, porque faz parte do projeto aprovado. E o escopo é
pequeno, porque é um domínio só.

O projeto de negócio vira o primeiro cliente do master data, e é ele que valida o
resultado na prática: se os vendedores reconhecem as próprias carteiras no painel, a
conciliação funcionou. Quando o próximo projeto de negócio chegar e precisar de clientes,
o domínio já está pronto, e ele só precisa puxar o domínio que falta, digamos, produtos. A
cada projeto, um domínio novo; a cada domínio, o próximo projeto fica mais barato (parte de
baixo da Figura~\ref{fig:master-data-abordagens}).

\begin{figure}[htbp]
  \centering
  % Master data: programa unico contra dominios puxados por projetos de negocio (cap. 13).
\begin{tikzpicture}[diagrama, x=1mm, y=1mm,
    titulo/.style={anchor=west, font=\footnotesize\scshape, text=cinza},
    linha/.style={anchor=west, font=\footnotesize, text=cinza}]

  % --- antes
  \node[titulo] at (0,40) {antes: um programa de master data};
  \node[caixa, anchor=west, minimum width=26mm] at (20,31) {Mapear fontes};
  \node[caixa, anchor=west, minimum width=34mm] at (46,31) {Conciliar nomes e IDs};
  \node[caixa, anchor=west, minimum width=20mm] at (80,31) {Governança};
  \node[destaquecaixa, anchor=west, minimum width=16mm] at (100,31) {Resultado};
  \draw[seta, draw=cinza] (20,24.5) -- (116,24.5);
  \node[rotulo, anchor=north] at (68,24) {meses ou anos até o primeiro resultado};

  % --- agora
  \node[titulo] at (0,12) {agora: puxado por um projeto de negócio};
  \node[linha] at (0,3) {Negócio};
  \node[linha] at (0,-9) {Master data};
  \node[caixa, anchor=west, minimum width=44mm] (pa) at (20,3) {Projeto de negócio A};
  \node[caixa, anchor=west, minimum width=46mm] (pb) at (70,3) {Projeto de negócio B};
  \node[destaquecaixa, anchor=west, minimum width=28mm] (cl) at (26,-9) {Clientes};
  \node[destaquecaixa, anchor=west, minimum width=28mm] (pr) at (80,-9) {Produtos};
  \draw[seta] (cl.north) -- (cl.north |- pa.south);
  \draw[seta] (pr.north) -- (pr.north |- pb.south);
  \draw[seta, dashed, draw=cinza] (cl.east) -- ($(pb.south west)+(4,0)$)
    node[rotulo, pos=0.45, below right] {reuso};
  \draw[seta, draw=cinza] (20,-16.5) -- (116,-16.5);
  \node[rotulo, anchor=north] at (68,-17)
    {um domínio por vez, cada um justificado por um projeto que precisa dele};
\end{tikzpicture}

  \caption{Dois jeitos de construir master data: um programa único, que entrega no
    fim, ou um domínio por vez, puxado por projetos de negócio que já precisam dele.}
  \label{fig:master-data-abordagens}
  \fonte{elaborado pelo autor.}
\end{figure}

A abordagem tem cuidados próprios. O primeiro é não deixar o prazo do projeto de negócio
atropelar a qualidade do domínio: um cadastro mestre feito às pressas, só para destravar
uma entrega, vira mais uma versão da verdade. O segundo é desenhar o domínio pensando no
reuso, e não só no primeiro projeto. O cadastro de clientes que nasce para o projeto
comercial vai ser usado pelo financeiro e pela logística, e precisa ter os atributos e as
regras que eles vão precisar. O terceiro é dar visibilidade ao projeto paralelo: ele entra
no quadro do time como um épico próprio, com histórias amarradas às do projeto de negócio
(\veja{sec:projeto-sustentacao}), para que o trabalho não fique escondido.

\rascunho{
  Relato completo de um caso real dessa abordagem: o projeto de negócio, o domínio que
  ele puxou, quanto tempo levou, quem validou e o que foi reaproveitado depois. É o relato
  mais importante do capítulo.
}

\section{O que a IA muda na conciliação}
\label{sec:master-data-ia}

A etapa que fazia o master data levar meses era a conciliação: decidir, registro por
registro, se dois cadastros representam a mesma entidade. É um problema com teoria antiga.
Em 1969, Ivan Fellegi e Alan Sunter formalizaram a conciliação probabilística de registros,
que compara os campos de cada par e calcula a probabilidade de os dois serem o mesmo
\parencite{fellegi1969}. As técnicas que vieram depois, de normalização, comparação e
bloqueio, estão reunidas em \textcite{christen2012}, e há bibliotecas de código aberto
que implementam esse modelo em escala \parencite{linacre2022}.

A IA acrescenta a essas técnicas a capacidade de lidar com o que as regras não pegam: nomes
abreviados, razões sociais escritas de jeitos diferentes, endereços incompletos, a filial
que aparece com o nome da fazenda. Ela pode sugerir pares prováveis, padronizar textos e,
principalmente, explicar por que dois registros parecem ser o mesmo, o que acelera muito
a revisão humana. A pessoa continua decidindo; a IA tira o trabalho braçal.

\rascunho{
  Validar com a minha experiência: em que medida a IA de fato encurtou a conciliação nos
  projetos de master data que eu acompanhei, e onde ela errou.
}

O efeito sobre a abordagem por domínio é direto. Conciliar um domínio inteiro dentro do
prazo de um projeto de negócio ficou viável, o que antes exigiria um time dedicado por
meses. Um motor simples de conciliação, com regras, similaridade e sugestões da IA, está
descrito no capítulo~\ref{cap:aplicacao-propria} (\veja{sec:motor-conciliacao}).

Há um cuidado que não pode ser esquecido. Cadastro de cliente tem dado pessoal, e o que vai
para um modelo de IA externo sai do ambiente da empresa. A limpeza do que pode e do que não
pode ser enviado precisa acontecer antes (\veja{sec:seguranca}).

\section{Governança mínima que funciona}
\label{sec:governanca-master-data}

Um domínio de master data sem governança volta a se degradar em poucos meses: alguém cria
um cliente novo direto no sistema de origem, outro corrige o nome na planilha, e a
correspondência deixa de valer. A governança não precisa ser grande, mas precisa existir
desde o primeiro domínio.

\begin{itemize}
  \item \textbf{Um dono de negócio por domínio}, que decide as regras: o que é um cliente
    ativo, quando dois cadastros são o mesmo, quais atributos são obrigatórios.
  \item \textbf{Regras de criação e alteração}: quem pode criar um cliente novo e em qual
    sistema, e como essa criação chega ao cadastro mestre.
  \item \textbf{Testes automáticos} de unicidade e de correspondência, rodando a cada carga
    (\veja{sec:tdd}).
  \item \textbf{A tabela de correspondência versionada e auditável}, num banco de dados e
    não numa planilha, com o registro de quem decidiu cada vínculo.
  \item \textbf{Crescimento junto com os domínios}: governança do tamanho do que já foi
    entregue, ampliada a cada domínio novo.
\end{itemize}

A governança completa, com o papel do data steward na era da IA, é o assunto do próximo
capítulo. Como construir uma aplicação própria para tudo isso, do motor de conciliação à
fila de revisão, fica para o capítulo~\ref{cap:aplicacao-propria}.

\begin{leituras}
  \leitura{dama2017}{O guia de referência de gestão de dados. O capítulo sobre dados
    mestres e de referência dá o vocabulário e as responsabilidades.}
  \leitura{loshin2009}{Um livro inteiro sobre master data, do caso de negócio à
    arquitetura e à governança.}
  \leitura{christen2012}{Para quem vai implementar a conciliação: técnicas de
    comparação, bloqueio e avaliação da qualidade dos pares encontrados.}
  \leitura{fellegi1969}{O artigo clássico que formalizou a conciliação probabilística
    de registros, ainda base de muitas ferramentas.}
\end{leituras}
